% Nguồn SGK Toán 9 Tập một — Luyện tập chung sau Bài 2, trang 19--20:
% https://cdn3.olm.vn/upload/thuvienso/4491668113-page-19-1773022927071.png
% https://cdn3.olm.vn/upload/thuvienso/4491668113-page-20-1773022927361.png

\section*{Luyện tập chung}

\begin{lesson}
    \textbf{Ví dụ 1}

    Giải hệ phương trình
    \[
        \begin{cases}
            0{,}5x+0{,}6y=0{,}4\\
            0{,}4x-0{,}9y=1{,}7.
        \end{cases}
    \]

    \textbf{Giải}

    Nhân hai vế của mỗi phương trình với 10, ta được:
    \[
        \begin{cases}
            5x+6y=4\\
            4x-9y=17.
        \end{cases}
        \tag{1}
    \]
    Ta giải hệ (1). Nhân hai vế của phương trình thứ nhất với 3 và nhân hai vế của phương trình thứ hai với 2, ta được hệ:
    \[
        \begin{cases}
            15x+18y=12\\
            8x-18y=34.
        \end{cases}
        \tag{2}
    \]
    Cộng từng vế hai phương trình của hệ (2) ta được $23x=46$, suy ra $x=2$.

    Thế $x=2$ vào phương trình thứ nhất của hệ (1), ta được $5\cdot2+6y=4$ hay $6y=-6$, suy ra $y=-1$.

    Vậy hệ phương trình đã cho có nghiệm là $(2;-1)$.

    \textbf{Ví dụ 2}

    Tìm các hệ số $x$, $y$ trong phản ứng hoá học đã được cân bằng sau:
    \[
        3\mathrm{Fe}+x\mathrm{O}_2\rightarrow y\mathrm{Fe}_3\mathrm{O}_4.
    \]

    \textbf{Giải}

    Vì số nguyên tử của Fe và O ở cả hai vế của phương trình phản ứng phải bằng nhau nên ta có hệ phương trình
    \[
        \begin{cases}
            3=3y\\
            2x=4y
        \end{cases}
        \quad\text{hay}\quad
        \begin{cases}
            y=1\\
            x=2y.
        \end{cases}
    \]
    Giải hệ này ta được $x=2$, $y=1$.

    \textbf{Ví dụ 3}

    Tìm hai số $a$ và $b$ để đường thẳng $y=ax+b$ đi qua hai điểm $A(-2;-1)$ và $B(2;3)$.

    \textbf{Giải}

    Đường thẳng $y=ax+b$ đi qua điểm $A(-2;-1)$ nên $a(-2)+b=-1$ hay $-2a+b=-1$.

    Tương tự, đường thẳng $y=ax+b$ đi qua điểm $B(2;3)$ nên $a\cdot2+b=3$ hay $2a+b=3$.

    Từ đó, ta có hệ phương trình với hai ẩn là $a$ và $b$:
    \[
        \begin{cases}
            -2a+b=-1\\
            2a+b=3.
        \end{cases}
    \]
    Cộng từng vế hai phương trình của hệ, ta được $2b=2$, suy ra $b=1$.

    Thay $b=1$ vào phương trình thứ nhất, ta có $-2a+1=-1$, suy ra $a=1$.

    Vậy với $a=1$, $b=1$ thì đường thẳng $y=ax+b$ đi qua hai điểm $A$, $B$ đã cho.
\end{lesson}

\begin{exercises}
    \subsection{Bài tập}

    \begin{questions}[resume=SGK]
        \item Cho hai phương trình:
        \[
            \begin{aligned}
                -2x+5y&=7; &&\text{(1)}\\
                4x-3y&=7. &&\text{(2)}
            \end{aligned}
        \]
        Trong các cặp số $(2;0)$, $(1;-1)$, $(-1;1)$, $(-1;6)$, $(4;3)$ và $(-2;-5)$, cặp số nào là:
        \begin{questions}
            \item Nghiệm của phương trình (1)?

            \answerline

            \item Nghiệm của phương trình (2)?

            \answerline

            \item Nghiệm của phương trình (1) và phương trình (2)?

            \answerline
        \end{questions}

        \item Giải các hệ phương trình sau bằng phương pháp thế:
        \begin{questions}
            \item $\begin{cases}2x-y=1\\x-2y=-1\end{cases};$

            \answerline
            \answerline
            \answerline

            \item $\begin{cases}0{,}5x-0{,}5y=0{,}5\\1{,}2x-1{,}2y=1{,}2\end{cases};$

            \answerline
            \answerline
            \answerline

            \item $\begin{cases}x+3y=-2\\5x-4y=28\end{cases}.$

            \answerline
            \answerline
            \answerline
        \end{questions}

        \item Giải các hệ phương trình sau bằng phương pháp cộng đại số:
        \begin{questions}
            \item $\begin{cases}5x+7y=-1\\3x+2y=-5\end{cases};$

            \answerline
            \answerline
            \answerline

            \item $\begin{cases}2x-3y=11\\-0{,}8x+1{,}2y=1\end{cases};$

            \answerline
            \answerline
            \answerline

            \item $\begin{cases}4x-3y=6\\0{,}4x+0{,}2y=0{,}8\end{cases}.$

            \answerline
            \answerline
            \answerline
        \end{questions}

        \item Tìm các hệ số $x$, $y$ trong phản ứng hoá học đã được cân bằng sau:
        \[
            4\mathrm{Al}+x\mathrm{O}_2\rightarrow y\mathrm{Al}_2\mathrm{O}_3.
        \]

        \answerline
        \answerline

        \item Tìm $a$ và $b$ sao cho hệ phương trình
        \[
            \begin{cases}
                ax+by=1\\
                ax+(2-b)y=3
            \end{cases}
        \]
        có nghiệm là $(1;-2)$.

        \answerline
        \answerline
        \answerline
    \end{questions}
\end{exercises}
